\chapter{Sensitivity to composition, mass, and injection parameters}
\label{chap:parameter_sensitivity}

With the physical steam model varied in \Cref{chap:steam_sensitivity},
this chapter turns to the remaining physical assumptions that are hard
to pin down independently: the planetary water inventory, the mass
range over which the analysis is applied, the runaway-greenhouse
fraction $f_\mathrm{rgh}$, the runaway threshold
$S_\mathrm{thresh}$, and the radius-measurement precision
$\sigma_R$.  Each is varied in turn while all other parameters remain
at the AGNI baseline of \Cref{chap:results}.  Figures in this chapter
are produced by \texttt{Src/figures/ch09\_parameter\_sensitivity.py}.


\section{Dependence on Planetary Water Content}
\label{sec:composition_dependence}

The amplitude of the injected radius inflation depends non-linearly on the
water mass fraction $x_{\mathrm{H_2O}}$.  This section isolates the
composition dependence at fixed steam prescription
(\texttt{agni\_updated}), moving from the fiducial value used by
\textcite{Schlecker2024} through a controlled sweep in
$x_{\mathrm{H_2O}}$ and into the water-rich regime accessible only to the
Aguichine and Pierrehumbert prescriptions.


\subsection{Fiducial water-to-rock ratio}
\label{ssec:wrr_fiducial}

% TODO --- scientific question:
%   Reference result: Δln Z when every planet is assigned x_H2O = 0.005,
%   the Schlecker+2024 fiducial value (about Earth's total water content).
% TODO --- plots / tables:
%   Not a new figure; refer back to the fiducial P4/P5 result from
%   \Cref{sec:hzied_plato} and quote its number here as the anchor for
%   the sweep in the next subsection.
% TODO --- pipeline call / notebook cell:
%   run_hzied_analysis(wrr=0.005, radius_inflation='agni_updated', ...)


\subsection{Variation of the water-to-rock ratio}
\label{ssec:wrr_sweep}

Holding the steam prescription (\texttt{agni\_updated}), the runaway-greenhouse
threshold ($S_\mathrm{thresh} = 280\,\mathrm{W\,m^{-2}}$), and the runaway fraction
($f_\mathrm{rgh} = 0.8$) fixed, we sweep the injected water-to-rock ratio
$x_{\mathrm{H_2O}}$ over every value supported by the AGNI+PALEOS ratio-method
steam table.  All remaining pipeline parameters are held at the P4/P5 baseline
configuration summarised in \Cref{tab:wrr_sweep_params}.

\begin{table}[H]
    \centering
    \caption[Fixed parameters of the $x_\mathrm{H_2O}$ sweep]{%
        Fixed configuration of the \texttt{agni\_updated} $x_\mathrm{H_2O}$
        sweep.  The only quantity varied across the sweep is
        $x_\mathrm{H_2O}$; every other parameter matches the P4/P5 baseline
        (\Cref{sec:hzied_plato}).}
    \label{tab:wrr_sweep_params}
    \small
    \begin{tabular}{lll}
        \toprule
        Parameter & Value & Note \\
        \midrule
        \multicolumn{3}{l}{\textit{Injection / model}} \\
        \texttt{radius\_inflation} & \texttt{agni\_updated} & AGNI+PALEOS steam table, ratio method \\
        \texttt{use\_dl21}         & \texttt{True}          & Dorn \& Lichtenberg (2021) water-incorporation \\
        $S_\mathrm{thresh}$        & $280\,\mathrm{W\,m^{-2}}$ & rocky HZ inner edge \\
        $f_\mathrm{rgh}$           & $0.8$                  & runaway greenhouse fraction \\
        \multicolumn{3}{l}{\textit{Observable cuts}} \\
        $M$                        & $[0.5,\,20]\,\Mearth$  & \texttt{agni\_updated} certified range \\
        depth                      & $\geq 60\,\mathrm{ppm}$ & relaxed cut (see \Cref{ssec:relaxed_depth}) \\
        \texttt{age\_mode}         & \texttt{sampled}       & Gaia FLAME + bootstrap fill \\
        \multicolumn{3}{l}{\textit{Inference}} \\
        window $w$                 & $25$                   & rolling SMA \\
        $\sigma_R / R$             & $4\%$ (P4), $3\%$ (P5) & PLATO Red Book \\
        $\sigma_S / S$             & $5\%$                  & Schlecker et al.\ (2024) \\
        \texttt{sma\_x\_mode}      & \texttt{standard}      & SMA-mean $S$ per window \\
        \texttt{sigma\_sma\_mode}  & \texttt{empirical}     & $\mathrm{std}(R_\mathrm{obs})/\sqrt{w}$ \\
        \texttt{nlive}             & $100$                  & dynesty live points \\
        seed base                  & $42$                   & 10 trials per grid point (seeds $42 + 13\,i$, $i=0\ldots9$) \\
        \bottomrule
    \end{tabular}
\end{table}

\Cref{fig:plato_wrr_sweep_lnk} shows the resulting Bayes factor $\ln\Delta Z$
as a function of $x_{\mathrm{H_2O}}$ for the two PLATO samples.  Each point
is the median across 10 independent noise realisations, and the shaded band
gives the 16--84\,\% inter-percentile range; the width of the band therefore
captures the combined jitter from measurement-noise draws and dynesty runs
at $n_\mathrm{live}=100$.  In practice the scatter is substantial: the full
range across 10 seeds spans ${\approx}\,20\text{--}30\,\ln\Delta Z$ units
on both samples, comparable to the systematic gain of moving from the
fiducial to the wettest grid value.  Any single-seed reading is therefore
best interpreted as one draw from this distribution rather than as a
definitive number.

Above the noise band, the trend is unambiguous: on both samples the median
evidence rises monotonically once $x_{\mathrm{H_2O}} \gtrsim 5\times10^{-3}$,
because the per-planet radius increment $\Delta R(x_{\mathrm{H_2O}})$
tabulated in the AGNI grid is roughly log-linear in $x_{\mathrm{H_2O}}$
up to the wet limit of the ratio method; a larger step relative to
$\sigma_R = 4\%$ (P4) or $3\%$ (P5) translates directly into a larger
Bayes factor at fixed planet count.  On P4 the median climbs from
$\ln\Delta Z \approx 49$ at the fiducial to ${\approx}\,93$ at
$4\times10^{-2}$; on P5 from ${\approx}\,18$ to ${\approx}\,49$ over the
same range.  The absolute gain is comparable on the two samples, but the
fractional gain is smaller on P5 because its larger sample already
saturates the information content of the fiducial step.

Below the fiducial value, the near-identical medians on both P4
($49.0$ at $10^{-4}$ vs $49.2$ at $5\times10^{-3}$) and P5 ($17.1$ vs
$18.4$) imply that, at the noise level of PLATO's observable samples,
the HZIED test is essentially insensitive to $x_{\mathrm{H_2O}}$ for any
Earth-like or drier composition: all such scenarios produce the same
statistically decisive discontinuity within the seed-to-seed scatter.
Only strongly water-enriched scenarios
($x_{\mathrm{H_2O}} \gtrsim 2\times10^{-2}$) leave a distinguishable
imprint on the Bayes factor once the noise band is accounted for.

\paragraph{Recovery of $S_\mathrm{thresh}$.}
The right panel of \Cref{fig:plato_wrr_sweep_lnk} shows how the
posterior-median $\hat{S}_\mathrm{thresh}$ recovered by dynesty tracks the
$\ln\Delta Z$ trend on the left.  At $x_{\mathrm{H_2O}} = 10^{-4}$ the
recovered value undershoots the injected $280\,\mathrm{W\,m^{-2}}$ by
${\approx}\,35\,\mathrm{W\,m^{-2}}$ on P4 ($\hat{S}_\mathrm{thresh}=246$)
and by ${\approx}\,23\,\mathrm{W\,m^{-2}}$ on P5 ($\hat{S}_\mathrm{thresh}=257$),
because the tiny radius step at that water content is not well localised
in $S$ and the posterior is pulled toward denser hot-side data at lower
$S_\mathrm{thresh}$.  As $x_{\mathrm{H_2O}}$ grows the discontinuity
sharpens and the recovery migrates back toward the injected value,
reaching $264$ (P4) and $272$ (P5) at $4\times10^{-2}$.  The convergence
of both $\hat{S}_\mathrm{thresh}$ and $\ln\Delta Z$ toward their
large-$x_{\mathrm{H_2O}}$ asymptotes is the same physical statement:
a larger radius step is both easier to detect (higher $\ln\Delta Z$)
and easier to localise (better $\hat{S}_\mathrm{thresh}$ recovery).
The seed-to-seed uncertainty on $\hat{S}_\mathrm{thresh}$ is
${\approx}\,20\,\mathrm{W\,m^{-2}}$ (16--84\,\%) across the grid,
smaller on P5 where the larger planet count reduces the posterior width.

\begin{figure}[H]
    \centering
    \includegraphics[width=\textwidth]{plato_wrr_sweep_lnk}
    \caption[$\ln\Delta Z$ and recovered $S_\mathrm{thresh}$ vs $x_\mathrm{H_2O}$]{%
        Left: Bayes factor $\ln\Delta Z$ as a function of the injected
        water-to-rock ratio $x_\mathrm{H_2O}$ for P4 (blue) and P5 (red).
        Right: the corresponding posterior-median of the recovered
        HZ-inner-edge threshold $\hat{S}_\mathrm{thresh}$; the dotted
        horizontal line is the injected value $280\,\mathrm{W\,m^{-2}}$.
        In both panels each marker is the median of 10 independent noise
        realisations, and the shaded band spans the 16th--84th percentile
        of those 10 draws.  All other pipeline parameters are as listed
        in \Cref{tab:wrr_sweep_params}; the dashed vertical line marks the
        Schlecker et al.\ (2024) fiducial $x_\mathrm{H_2O} = 5\times10^{-3}$.
        The recovered $\hat{S}_\mathrm{thresh}$ systematically undershoots
        the injected value at small $x_\mathrm{H_2O}$ and converges toward it
        as the radius step grows, tracking the same trend as $\ln\Delta Z$.}
    \label{fig:plato_wrr_sweep_lnk}
\end{figure}

\Cref{fig:plato_wrr_sweep_panels} shows the detected planet population for the
two extremes of the grid and the fiducial value.  The vertical grey lines
mark the steam-atmosphere inflation applied per planet: they lengthen
visibly from $x_\mathrm{H_2O}=10^{-4}$ (nearly imperceptible) through the
fiducial $5\times10^{-3}$ to $4\times10^{-2}$, where the hot-side
inflation grows to $\gtrsim 0.15\,\Rearth$ per planet and produces a
clearly visible discontinuity at $S_\mathrm{thresh}=280\,\mathrm{W\,m^{-2}}$.

\begin{figure}[H]
    \centering
    \includegraphics[width=\textwidth]{plato_wrr_sweep_panels}
    \caption[Detected planet sample for key $x_\mathrm{H_2O}$ values]{%
        Detected planet sample for P4 (top row) and P5 (bottom row) at
        $x_\mathrm{H_2O} \in \{10^{-4},\,5\times10^{-3},\,4\times10^{-2}\}$
        (left to right).  Vertical grey lines connect each inflated planet's
        pre-inflation radius to its observed radius, illustrating how the
        steam-atmosphere increment grows with $x_\mathrm{H_2O}$.  The
        recovered $\ln\Delta Z$ is annotated in each panel title.  All other
        parameters as in \Cref{tab:wrr_sweep_params}.}
    \label{fig:plato_wrr_sweep_panels}
\end{figure}

% Figures produced by:
%   Src/figures/plato_wrr_sweep.py
%     → Figures/generated/plato_wrr_sweep_lnk.pdf
%     → Figures/generated/plato_wrr_sweep_panels.pdf


\subsection{Water-rich compositions}
\label{ssec:water_rich}

% TODO --- scientific question:
%   In the water-rich regime accessible only to the Aguichine and
%   Pierrehumbert prescriptions, does the discontinuity strengthen,
%   weaken, or fundamentally change shape?
% TODO --- plots / tables:
%   Δln Z as a function of x_H2O (or WMF) in the range 0.05 -- 0.5, using
%   radius_inflation='aguichine'. Comparison plot juxtaposing the Aguichine
%   result at WMF = 0.5 with the Pierrehumbert Hycean case.
% TODO --- pipeline call / notebook cell:
%   run_hzied_analysis(wrr=<large>, radius_inflation='aguichine', ...)
%   run_hzied_analysis(radius_inflation='hycean_pierrehumbert', ...)


% =======================================================================
\section{Dependence on Planet Mass and Applicability Limits}
\label{sec:mass_dependence}

The steam-atmosphere prescriptions differ substantially in the planet
mass range over which they are certified.  This section separates two
distinct effects: (i) a \emph{physical} change in the recovered
$\Delta \ln Z$ due to the mass dependence of the radius inflation itself,
and (ii) a \emph{sample-selection} change due to the observable mass cut
$M_p \in [M_{\min}, M_{\max}]$ that must accompany each prescription.

% TODO --- scientific question:
%   How much of the model-to-model spread in Δln Z shown in
%   \Cref{sec:steam_model} is driven by the physics of inflation as a
%   function of planet mass, versus by the differing observable mass
%   ranges enforced by the certified grid limits of each prescription?
%
% TODO --- plots / tables:
%   (i)  ΔR(M) at fixed water mass fraction for each prescription,
%        showing the mass range and the shape of the inflation curve.
%   (ii) Δln Z on P4 as the upper mass cut M_max is varied
%        (e.g. 2, 3, 5, 10, 20 M_earth) within a single prescription
%        (agni_updated). Highlights the sample-selection effect at fixed
%        physics.
%   (iii) Table: N_planets and Δln Z for each (prescription, M_max)
%        combination, distinguishing the sample-selection effect from the
%        physical-inflation effect.
%
% TODO --- pipeline call / notebook cell:
%   run_hzied_analysis(sample='P4', radius_inflation='agni_updated',
%                      M_min=0.5, M_max=<value>, ...)


% =======================================================================
\section{Other Key Model Parameters}
\label{sec:other_parameters}

The remaining model parameters that control the injected signal are
grouped into physically interpretable blocks.  Each is isolated by
holding the other axes at their fiducial values.


\subsection{Runaway fraction \texorpdfstring{$f_{\mathrm{rgh}}$}{f\_rgh}}
\label{ssec:f_rgh}

The runaway-greenhouse fraction $f_\mathrm{rgh}$ absorbs the population's
unknown steam-atmosphere lifetime: it is the fraction of hot-side
planets ($S_\mathrm{abs} > S_\mathrm{thresh}$) that are currently in an
inflated runaway state.  Physically, $f_\mathrm{rgh}$ interpolates
between the two extreme scenarios that H1 and H0 encode: at
$f_\mathrm{rgh}=0$ no hot-side planet is inflated and the H1 step
degenerates to the flat H0 model (so $\ln\Delta Z$ collapses to sampler
noise around zero); at $f_\mathrm{rgh}=1$ every hot-side planet is
inflated and the injected step reaches its maximum amplitude
$\Delta R = f_{dR}(x_{\mathrm{H_2O}})$.  For intermediate values the
average injected step scales linearly, $\Delta R_\mathrm{eff}
= f_\mathrm{rgh} \cdot f_{dR}(x_{\mathrm{H_2O}})$, so the expected
scaling of the Bayes factor is $\ln\Delta Z \propto f_\mathrm{rgh}^2$
at fixed noise (since evidence for a Gaussian mean shift scales with
the squared amplitude relative to the error).

We sweep $f_\mathrm{rgh} \in \{0.2,\,0.4,\,0.6,\,0.8,\,1.0\}$ at the
otherwise-fiducial \texttt{agni\_updated} configuration
(\Cref{tab:frgh_sweep_params}); $f_\mathrm{rgh}=0$ is omitted from the
grid because it degenerates H1 into H0 and contributes only sampler
noise, which would compress the plot without adding information.  The
Schlecker et al.\ (2024) fiducial value is $0.8$.

\begin{table}[H]
    \centering
    \caption[Fixed parameters of the $f_\mathrm{rgh}$ sweep]{%
        Fixed configuration of the \texttt{agni\_updated} $f_\mathrm{rgh}$
        sweep.  Only $f_\mathrm{rgh}$ varies across the grid; every other
        parameter matches the P4/P5 baseline used in
        \Cref{ssec:wrr_sweep}.}
    \label{tab:frgh_sweep_params}
    \small
    \begin{tabular}{lll}
        \toprule
        Parameter & Value & Note \\
        \midrule
        \multicolumn{3}{l}{\textit{Injection / model}} \\
        \texttt{radius\_inflation} & \texttt{agni\_updated} & AGNI+PALEOS steam table \\
        \texttt{use\_dl21}         & \texttt{True}          & Dorn \& Lichtenberg (2021) water incorporation \\
        $x_\mathrm{H_2O}$          & $0.005$                & Schlecker+2024 fiducial \\
        $S_\mathrm{thresh}$        & $280\,\mathrm{W\,m^{-2}}$ & rocky HZ inner edge \\
        \multicolumn{3}{l}{\textit{Observable cuts}} \\
        $M$                        & $[0.5,\,20]\,\Mearth$  & \texttt{agni\_updated} certified range \\
        depth                      & $\geq 60\,\mathrm{ppm}$ & relaxed cut (\Cref{ssec:relaxed_depth}) \\
        \texttt{age\_mode}         & \texttt{sampled}       & FLAME + bootstrap fill \\
        \multicolumn{3}{l}{\textit{Inference}} \\
        window $w$                 & $25$                   & rolling SMA \\
        $\sigma_R / R$             & $4\%$ (P4), $3\%$ (P5) & PLATO Red Book \\
        $\sigma_S / S$             & $5\%$                  & Schlecker+2024 \\
        \texttt{sma\_x\_mode}      & \texttt{standard}      & mean $S$ per window \\
        \texttt{sigma\_sma\_mode}  & \texttt{empirical}     & $\mathrm{std}(R_\mathrm{obs})/\sqrt{w}$ \\
        \texttt{nlive}             & $100$                  & dynesty live points \\
        seed base                  & $42 + 200\,i$ ($i=0..4$) & unique per grid point; 10 trials each ($+13k$, $k=0..9$) \\
        \bottomrule
    \end{tabular}
\end{table}

\Cref{fig:plato_frgh_sweep_lnk} shows the recovered $\ln\Delta Z$
(left) and posterior-median $\hat{S}_\mathrm{thresh}$ (right) as
functions of $f_\mathrm{rgh}$ for the two PLATO samples.  Each marker
is the median across $10$ independent noise realisations; the shaded
band is the $16$--$84\,\%$ range, capturing measurement-noise scatter
plus dynesty variance at $n_\mathrm{live}=100$ (same construction as
\Cref{ssec:wrr_sweep}).  The multi-seed medians are collected in
\Cref{tab:frgh_sweep_results}.

\begin{table}[H]
    \centering
    \caption[Multi-seed medians of the $f_\mathrm{rgh}$ sweep]{%
        Multi-seed medians of the $f_\mathrm{rgh}$ sweep.  $\ln\Delta Z$
        columns quote the median across $10$ noise realisations with the
        per-seed min/max bracket; $\hat{S}_\mathrm{th}$ columns are the
        median of the posterior-median values across the same $10$ draws
        (W\,m$^{-2}$).  The injected threshold is $280\,\mathrm{W\,m^{-2}}$.
        Planet counts are $N_\mathrm{P4}=139$ and $N_\mathrm{P5}=266$
        across the grid (independent of $f_\mathrm{rgh}$).}
    \label{tab:frgh_sweep_results}
    \small
    \begin{tabular}{lcccc}
        \toprule
        $f_\mathrm{rgh}$ & P4 $\ln\Delta Z$ (med [min,max]) & P4 $\hat{S}_\mathrm{th}$
                         & P5 $\ln\Delta Z$ (med [min,max]) & P5 $\hat{S}_\mathrm{th}$ \\
        \midrule
        $0.20$ & $29.2\;[20.4,\,36.9]$ & $229$ & $14.9\;[\phantom{0}8.4,\,20.0]$ & $259$ \\
        $0.40$ & $33.3\;[24.8,\,45.2]$ & $238$ & $15.4\;[11.1,\,26.4]$           & $264$ \\
        $0.60$ & $39.4\;[25.0,\,49.0]$ & $229$ & $17.6\;[13.0,\,23.4]$           & $260$ \\
        $0.80$ & $39.0\;[32.5,\,50.6]$ & $248$ & $17.8\;[16.0,\,21.9]$           & $261$ \\
        $1.00$ & $50.5\;[42.1,\,76.5]$ & $243$ & $21.7\;[14.2,\,28.7]$           & $262$ \\
        \bottomrule
    \end{tabular}
\end{table}

Three concrete observations emerge:

\begin{enumerate}
    \item \textbf{$\ln\Delta Z$ grows sub-quadratically with
    $f_\mathrm{rgh}$.}  The naive expectation for a Gaussian mean shift
    scaling with the squared step amplitude,
    $\ln\Delta Z \propto f_\mathrm{rgh}^2$, would predict a factor
    ${\sim}25$ between $f_\mathrm{rgh}=0.2$ and $f_\mathrm{rgh}=1.0$.
    The observed medians grow only by a factor of ${\sim}1.7$--$1.9$
    (P4: $29\to 51$; P5: $15\to 22$).  Two reasons: (i) the SMA
    likelihood is not a pure two-mean test --- it operates on a
    rolling window and therefore integrates the step over several
    S bins, so the effective SNR grows more slowly than the raw step
    amplitude; (ii) even at $f_\mathrm{rgh}=0.2$ the hot side already
    contains enough inflated planets to be statistically distinct from
    the cold side, so the transition from "detectable" to "decisive"
    is already made at modest fractions.
    \item \textbf{Recovered $\hat{S}_\mathrm{thresh}$ is essentially
    flat in $f_\mathrm{rgh}$.}  On P4 it stays at
    $229$--$248\,\mathrm{W\,m^{-2}}$ across the grid, and on P5 at
    $259$--$264\,\mathrm{W\,m^{-2}}$, in both cases with a seed-to-seed
    scatter narrower than the mode-to-mode variation.  This is a
    reassuring orthogonality: the \emph{location} of the step is set by
    the planet distribution in $S_\mathrm{obs}$ and is unaffected by
    the \emph{amplitude} of the injected step.  Both samples continue
    to show the modest undershoot ($\hat{S} < 280\,\mathrm{W\,m^{-2}}$)
    seen throughout this chapter, driven by the higher density of hot-side
    planets that pulls the fitted transition toward smaller $S$.
    \item \textbf{The seed-to-seed band narrows on P5 more than on
    P4.}  The P5 band widths are $\lesssim 12\,\ln\Delta Z$-units
    across the grid, while the P4 band widths are typically
    $\sim 15$--$35$ units.  This reflects P5's larger planet count
    (${\sim}\!2\times$ that of P4) reducing the measurement-noise
    contribution to $\sigma_\mathrm{SMA}$.  Even in the fiducial
    scenario the two samples remain in different noise regimes,
    which motivates keeping them presented separately throughout the
    chapter.
\end{enumerate}

Taken together with the wrr sweep (\Cref{ssec:wrr_sweep}), we can
place order-of-magnitude constraints on the joint sensitivity of the
Bayes factor: doubling $f_\mathrm{rgh}$ at fixed $x_{\mathrm{H_2O}}$
buys roughly the same ${\sim}\!10$--$25$-unit gain in $\ln\Delta Z$
as raising $x_{\mathrm{H_2O}}$ by half an order of magnitude at fixed
$f_\mathrm{rgh}$.  Neither knob is dominant; both saturate softly
because the inference is already SMA-driven and window-limited.

\Cref{fig:plato_frgh_sweep_panels} shows the detected planet sample
for one representative trial at $f_\mathrm{rgh} \in \{0.2,\,0.8,\,1\}$:
the fraction of coloured "inflated" points and the total length of
the grey inflation lines both grow visibly with $f_\mathrm{rgh}$,
sharpening the step at $S_\mathrm{thresh}$.

\begin{figure}[H]
    \centering
    \includegraphics[width=\textwidth]{plato_frgh_sweep_lnk}
    \caption[$\ln\Delta Z$ and $\hat{S}_\mathrm{thresh}$ vs $f_\mathrm{rgh}$]{%
        Left: Bayes factor $\ln\Delta Z$ as a function of the injected
        runaway-greenhouse fraction $f_\mathrm{rgh}$ for P4 (blue) and
        P5 (red).  Right: the corresponding posterior-median of the
        recovered runaway-GH threshold $\hat{S}_\mathrm{thresh}$; the
        dotted horizontal line is the injected value
        $280\,\mathrm{W\,m^{-2}}$.  In both panels each marker is the
        median of $10$ independent noise realisations and the shaded
        band spans the $16$--$84\,\%$ range of those draws.  The
        dashed vertical line marks the Schlecker et al.\ (2024)
        fiducial $f_\mathrm{rgh}=0.8$.  All other pipeline parameters
        are as listed in \Cref{tab:frgh_sweep_params}.}
    \label{fig:plato_frgh_sweep_lnk}
\end{figure}

\begin{figure}[H]
    \centering
    \includegraphics[width=\textwidth]{plato_frgh_sweep_panels}
    \caption[Detected planet sample for key $f_\mathrm{rgh}$ values]{%
        Detected planet sample for P4 (top row) and P5 (bottom row) at
        $f_\mathrm{rgh} \in \{0.2,\,0.8,\,1\}$ (left to right).
        Vertical grey lines connect each inflated planet's rocky
        baseline radius to its observed radius; the fraction of hot-side
        planets carrying an inflation vector grows with $f_\mathrm{rgh}$,
        driving the recovered $\ln\Delta Z$.  All other parameters as
        in \Cref{tab:frgh_sweep_params}.}
    \label{fig:plato_frgh_sweep_panels}
\end{figure}

% Figures produced by:
%   Src/figures/plato_frgh_sweep.py
%     → Figures/generated/plato_frgh_sweep_lnk.pdf
%     → Figures/generated/plato_frgh_sweep_panels.pdf


\subsection{Runaway threshold \texorpdfstring{$S_{\mathrm{thresh}}$}{S\_thresh}}
\label{ssec:s_thresh}

% TODO --- scientific question:
%   How does varying the injected runaway threshold move the recovered
%   posterior, and how well can the inference recover the injected value?
% TODO --- plots / tables:
%   Recovered S_thresh posterior for injected values in {200, 250, 280,
%   310, 350} W/m^2. Injection--recovery diagonal.
% TODO --- pipeline call / notebook cell:
%   run_hzied_analysis(S_thresh=<value>, ...) — Generator auto-run for
%   non-fiducial thresholds.
%   → notebooks/hzied_pipeline.ipynb Test 4 (cell 12)


\subsection{Radius measurement precision \texorpdfstring{$\sigma_R$}{sigma\_R}}
\label{ssec:sigma_r}

% TODO --- scientific question:
%   How strongly does the detection depend on the per-planet radius
%   precision assumed for each PLATO sample (Red Book values: P1 1\%,
%   P2 0.8\%, P4 4\%, P5 3\%)?
% TODO --- plots / tables:
%   Δln Z on P4 (and P5) as sigma_R_frac is varied around the Red Book
%   default (e.g. 1\%, 2\%, 3\%, 4\%, 5\%).
% TODO --- pipeline call / notebook cell:
%   run_hzied_analysis(sigma_R_frac=<value>, ...) overriding the
%   per-sample default in _SIGMA_R.


\subsection{Observable-cut choices}
\label{ssec:cuts}

% TODO --- scientific question:
%   How sensitive is the detection to the choice of transit-depth
%   threshold and mass cuts used to define the observable subset?
% TODO --- plots / tables:
%   (i)  Planet counts and Δln Z as depth_min_ppm is varied (80 → 60 → 40).
%   (ii) Planet counts and Δln Z as M_min, M_max is varied at fixed
%        inflation model.
%
% TODO --- pipeline call / notebook cell:
%   run_hzied_analysis(depth_min_ppm=<value>, M_min=<value>,
%                      M_max=<value>, ...)
%   → notebooks/hzied_pipeline.ipynb transit-depth sweep (cell 26);
%     wp1_detection_limits.ipynb for the full depth exploration.


\subsection{Hypothesis-test methodology}
\label{ssec:methodology}

% TODO --- scientific question:
%   Is the recovered Δln Z robust to the choice of statistical
%   estimator (rolling vs independent bins, propagated vs empirical
%   sigma, cold_edge vs standard x-mode)?
% TODO --- plots / tables:
%   Small robustness table: recovered Δln Z under each combination.
% TODO --- pipeline call / notebook cell:
%   run_hzied_analysis(sigma_sma_mode=..., independent_bins=...,
%                      sma_x_mode=..., ...)
%   → notebooks/hzied_pipeline.ipynb Test 6 (cell 16) and cells 28--32.


